\section{Wiles 证明策略概述}
\label{sec:ch10-wiles-strategy}
% ============================================================================

\providecommand{\rhobar}{\bar{\rho}}
\providecommand{\Sel}{\operatorname{Sel}}

Wiles 证明最令人震撼的地方，
不只在于它解决了一个三百多年的猜想，
更在于它把“椭圆曲线模性”这样一句几何断言
彻底翻译成了一个关于 Galois 表示形变与 Hecke 代数的同构问题。
从外部看，
费马大定理似乎与 Hecke 局部化、Selmer 群、完整交环风马牛不相及；
但一旦进入 Wiles 的视角，
这些对象反而比原方程本身更接近答案。

\textbf{我们遇到了什么问题？}
模性定理要求从一条半稳定椭圆曲线 $E/\Q$ 出发，
证明它对应某个新形式 $f$。
直接构造 $f$ 几乎不可能，
因为模形式空间太大，
而椭圆曲线的算术数据又以局部 Galois 条件的形式出现。
Wiles 的突破在于：
既然第\ref{ch:galois}章已经告诉我们“模性”可以写成 Galois 表示同构，
那就不妨直接比较两类表示的 \emph{提升空间}。

\textbf{核心思想是什么？}
先固定一个 residual 表示
\[
  \rhobar\colon G_\Q\to \GL_2(\F_\ell).
\]
一边，考虑所有满足给定局部条件的 $\ell$-进提升，
它们由某个泛形变环 $R$ 参数化；
另一边，考虑所有满足同样局部条件的模形式，
它们在局部 Hecke 代数中编码成一个环 $T$。
如果能证明
\[
  R\cong T,
\]
那么每一个允许的 Galois 提升都自动是模的。
Wiles 的核心定理，
不是“直接找到了对应的新形式”，
而是证明这两个看似来自不同世界的局部环其实同构。

\textbf{引入之后能得到什么？}
一旦 $R=T$，
椭圆曲线的 $\ell$-进表示落在 $R$ 上的一个点就对应于 $T$ 上的一个 Hecke 特征值系统，
也就对应于一个新形式。
这把模性问题从“构造对象”变成了“证明两个变形空间一样大”。
接下来，Langlands--Tunnell 提供起点，
Taylor--Wiles 提供比较环大小的方法，
而 $3$--$5$ 技巧则处理最初 residual 表示可能可约的例外情形。

\begin{figure}[ht]
\centering
\begin{tikzpicture}[>=Stealth, font=\small, node distance=1.2cm]
  \node[draw, rounded corners, fill=orange!14, minimum width=4.8cm, minimum height=0.9cm] (res) {残余表示 $\rhobar_{E,3}$};
  \node[draw, rounded corners, fill=yellow!18, below=of res, minimum width=4.8cm, minimum height=0.9cm] (lt) {Langlands--Tunnell 给出残余模性};
  \node[draw, rounded corners, fill=green!10, below left=1.3cm and 1.8cm of lt, minimum width=3.8cm, minimum height=0.9cm] (r) {形变环 $R$};
  \node[draw, rounded corners, fill=blue!8, below right=1.3cm and 1.8cm of lt, minimum width=3.8cm, minimum height=0.9cm] (t) {Hecke 环 $T$};
  \node[draw, rounded corners, fill=purple!10, below=2.0cm of lt, minimum width=4.2cm, minimum height=0.9cm] (rt) {Taylor--Wiles 证明 $R=T$};
  \node[draw, rounded corners, fill=red!8, below=of rt, minimum width=4.6cm, minimum height=0.9cm] (mod) {所有允许提升都来自模形式};
  \node[draw, rounded corners, fill=gray!14, below=of mod, minimum width=4.8cm, minimum height=0.9cm] (trick) {$3$--$5$ 技巧补足可约例外};

  \draw[->, very thick] (res) -- (lt);
  \draw[->, very thick] (lt) -- (r);
  \draw[->, very thick] (lt) -- (t);
  \draw[->, very thick] (r) -- (rt);
  \draw[->, very thick] (t) -- (rt);
  \draw[->, very thick] (rt) -- (mod);
  \draw[->, very thick] (mod) -- (trick);
\end{tikzpicture}
\caption{Wiles 证明策略的逻辑骨架：先有残余模性，再比较形变环与 Hecke 环，最后用 $R=T$ 把每个允许的提升都判定为模。}
\label{fig:ch10-wiles-flow}
\end{figure}


现在按逻辑顺序概述这条证明链。

\subsection*{第一步：先从一个 residual 表示起步}
对半稳定椭圆曲线 $E/\Q$，取 $\ell=3$，得到
\[
  \rhobar_{E,3}\colon G_\Q\to \GL_2(\F_3).
\]
若它绝对不可约，
则其射影像落在 $\PGL_2(\F_3)\cong S_4$ 中，
从而正好落入 Langlands--Tunnell 定理可处理的“可解射影像”范围。
这个定理给出 Wiles 证明真正的起跑线：
先证明残余表示已经是模的，
然后再讨论它的所有 $3$-进提升。

\begin{lemma}[Langlands--Tunnell]
\label{lem:ch10-langlands-tunnell}
设
\[
  \rhobar\colon G_\Q\to \GL_2(\F_3)
\]
是奇的、绝对不可约的连续表示，且其射影像
\[
  \mathbb P\rhobar(G_\Q)\subset \PGL_2(\F_3)
\]
是可解群。
则 $\rhobar$ 是模的：存在一个权 $1$ 的归一化特征 cusp form $g$ 与 $g$ 的某个 $3$ 进素点 $\lambda$，使得
\[
  \rhobar_{g,\lambda}\cong \rhobar.
\]
把 $g$ 乘以 Hasse 不变量后，也可把这同一个残余表示看成来自某个权 $2$ 模形式。
\end{lemma}

\begin{proof}
直觉上，这个定理说的是：当一个二维 Artin 表示的射影像还停留在可解世界里时，
Langlands 的可解基变换与 Tunnell 对八面体情形的补足，
足以把它送到自守一侧。
下面把这句话拆成四步。

\textbf{第一步：射影像确实可解。}
群 $\PGL_2(\F_3)$ 的阶为
\[
  \frac{(3^2-1)(3^2-3)}{3-1}=24,
\]
而它与 $S_4$ 同构。
另一方面，$S_4$ 有正规列
\[
  S_4\triangleright A_4\triangleright V_4\triangleright 1,
\]
其中
\[
  S_4/A_4\cong C_2,
  \qquad
  A_4/V_4\cong C_3,
  \qquad
  V_4/1\cong C_2\times C_2.
\]
这些商群全是 abelian，故 $S_4$ 可解，因而任何射影像子群也可解。

\textbf{第二步：把 mod $3$ 表示提升成复 Artin 表示。}
设 $G=\rhobar(G_\Q)$，则 $G$ 是有限群。
因为 $\rhobar$ 绝对不可约，$G$ 的二维表示在 $\overline{\F}_3$ 上不可约。
对有限群表示应用 Brauer 提升，可得一个二维复表示
\[
  r\colon G_\Q\twoheadrightarrow G\hookrightarrow \GL_2(\C),
\]
其特征标在与 $3$ 互素的元素上是 $\rhobar$ 特征标的 Teichm\"uller 提升。
特别地，$r$ 仍是不可约的，且因为复共轭 $c$ 满足
\[
  \det \rhobar(c)=-1,
\]
故也有 $\det r(c)=-1$；于是 $r$ 是奇的二维 Artin 表示。
它的射影像与 $\rhobar$ 的射影像相同，因此也是可解的。

\textbf{第三步：应用 Langlands--Tunnell 的可解 Artin 模性。}
Langlands 证明了可解扩张上的基变换与若干可解像情形的函子性；
Tunnell 进一步补足了射影像为八面体型的最后情形。
合起来得到：每个奇的、不可约的二维复 Artin 表示，只要其射影像可解，
就对应一个权 $1$ 的 cuspidal 自守表示，亦即一个权 $1$ 的归一化特征模形式 $g$。
Deligne--Serre 把这一定理表述为
\[
  \rho_g\cong r,
\]
且 $g$ 的级别恰等于 $r$ 的 Artin 导子。

\textbf{第四步：模 $3$ 约化回到原来的 residual 表示。}
取 $g$ 的系数域中一个位于 $3$ 上方的素点 $\lambda$。
因为 $\rho_g$ 与 $r$ 同构，而 $r$ 正是 $\rhobar$ 的 Brauer 提升，
故对几乎所有 Frobenius 元都有
\[
  \operatorname{tr}(\rho_g(\Frob_p))\equiv \operatorname{tr}(\rhobar(\Frob_p))\pmod{\lambda},
\]
且行列式也同余。
由 Chebotarev 密度定理与半单化唯一性，得到
\[
  \rhobar_{g,\lambda}^{\mathrm{ss}}\cong \rhobar^{\mathrm{ss}}.
\]
由于 $\rhobar$ 已绝对不可约，半单化可去掉，于是
\[
  \rhobar_{g,\lambda}\cong \rhobar.
\]
这就证明了 $\rhobar$ 的模性。

把这一步用于 $\rhobar_{E,3}$ 时，奇性来自 Weil 配对给出的
\[
  \det \rhobar_{E,3}=\chi_3,
\]
而若 $\rhobar_{E,3}$ 绝对不可约，则其射影像自然是 $S_4$ 的某个子群。
因此 Langlands--Tunnell 给出一个模形式 $f_0$，满足
\[
  \rhobar_{f_0,3}\cong \rhobar_{E,3},
\]
这正是 Wiles 提升定理的起步点。
\end{proof}

\subsection*{第二步：把“所有可能提升”做成一个环}
固定 $\rhobar_{E,3}$ 之后，
我们不只看某一条提升，
而是看所有满足局部条件的提升
\[
  \rho\colon G_\Q\to \GL_2(A),
\]
其中 $A$ 是完备局部环。
这些提升的同构类组成一个形变函子；
Mazur 证明，在温和不可约性条件下它由一个泛环 $R$ 表示。
因此，$R$ 的每一个局部点都对应一条“可能来自椭圆曲线或模形式”的 $3$-进表示。

\subsection*{第三步：模形式一侧出现 Hecke 环 $T$}
把满足相同局部条件的权 $2$ 模形式收集起来，
让 Hecke 算子在对应空间上作用，
再在与 $\rhobar_{E,3}$ 相容的极大理想处局部化并完备化，
得到 Hecke 环 $T$。
第\ref{ch:hecke}章已经告诉我们，
特征形式由 Hecke 特征值系统决定；
因此 $T$ 正是“模的提升”在局部环层面的编码。
由 universal property，存在自然满同态
\[
  R\twoheadrightarrow T.
\]
于是问题变成：这是否其实是同构？

\subsection*{第四步：数值判据与 Taylor--Wiles 打补丁}
Wiles 的关键观察是：
既然已经有自然满射 $R\twoheadrightarrow T$，
那就不必直接写出两个环的全部方程；
只要证明它们在“切空间大小”和“模形式同余大小”这两个数值上完全吻合，
满射就会被迫成为同构。

\begin{theorem}[Wiles 数值判据]
\label{thm:ch10-numerical-criterion}
设
\[
  \phi\colon R\twoheadrightarrow T
\]
是完备局部 $\Z_\ell$-代数的满射，
并设 $\pi\colon T\to \Z_\ell$ 是一个增广，$\eta=\pi\circ\phi\colon R\to \Z_\ell$。
记
\[
  I_R=\ker\eta,
  \qquad
  I_T=\ker\pi.
\]
定义
\[
  \Phi_T=\pi\bigl(\operatorname{Ann}_T(I_T)\bigr),
  \qquad
  \Phi_R=\eta\bigl(\operatorname{Ann}_R(I_R)\bigr),
\]
它们分别称为 $T$ 与 $R$ 的 congruence ideal。
在 Wiles 所处理的一维 Cohen--Macaulay 情形中，有等价命题：
\[
  \phi\ \text{是同构}
  \iff
  \#(\Z_\ell/\Phi_T)\le \#(\Z_\ell/\Phi_R).
\]
等价地，
\[
  \operatorname{length}_{\Z_\ell}(\Z_\ell/\Phi_T)
  \le
  \operatorname{length}_{\Z_\ell}(\Z_\ell/\Phi_R).
\]
若进一步 $T$ 是 complete intersection，
则上述不等式自动成为等号，且 $R$ 也是 complete intersection，因此 $R\cong T$。
\end{theorem}

\begin{proof}
直觉上，$\Phi_T$ 衡量的是：给定增广点以后，模形式侧还能与多少别的 Hecke 特征值发生同余；
而 $\Phi_R$ 衡量的是：形变侧在该点附近还保留多少无穷小自由度。
数值判据的内容就是：若模形式侧并不比形变侧“更小”，
那么满射 $R\twoheadrightarrow T$ 就已经没有核了。

先看最容易的一向。
若 $\phi$ 是同构，
则 $I_R\cong I_T$，并且
\[
  \operatorname{Ann}_R(I_R)\cong \operatorname{Ann}_T(I_T).
\]
经由增广映到 $\Z_\ell$ 后立刻得到
\[
  \Phi_R=\Phi_T,
\]
故两边长度相等。

下面设
\[
  \operatorname{length}_{\Z_\ell}(\Z_\ell/\Phi_T)
  \le
  \operatorname{length}_{\Z_\ell}(\Z_\ell/\Phi_R).
\]
记 $K=\ker\phi$。
因为 $\phi$ 与增广相容，$\phi(I_R)=I_T$，于是有正合列
\[
  K/(I_RK)\longrightarrow I_R/I_R^2\longrightarrow I_T/I_T^2\longrightarrow 0.
\]
因此 $I_R/I_R^2\twoheadrightarrow I_T/I_T^2$，从而必有
\[
  \operatorname{length}_{\Z_\ell}(I_T/I_T^2)
  \le
  \operatorname{length}_{\Z_\ell}(I_R/I_R^2).
\]
在 Wiles 的交换代数引理中，congruence ideal 正是这些余切模的 Fitting ideal：
\[
  \Phi_A=\operatorname{Fitt}_{\Z_\ell}(I_A/I_A^2)
  \qquad (A=R,T),
\]
故对一维 Cohen--Macaulay 局部环，
\[
  \operatorname{length}_{\Z_\ell}(\Z_\ell/\Phi_A)
  =
  \operatorname{length}_{\Z_\ell}(I_A/I_A^2).
\]
把它代入假设，得到
\[
  \operatorname{length}_{\Z_\ell}(I_T/I_T^2)
  =
  \operatorname{length}_{\Z_\ell}(I_R/I_R^2).
\]
既然前面已有一个满射
\[
  I_R/I_R^2\twoheadrightarrow I_T/I_T^2,
\]
而两边长度相等，这个满射只能是同构。
于是
\[
  K/(I_RK)=0.
\]
由于 $I_R$ 包含在局部环 $R$ 的 Jacobson 根中，Nakayama 引理推出 $K=0$。
因此 $\phi$ 单且满，也就是同构。

最后讨论 complete intersection 情形。
若 $T$ 是 complete intersection，
则其余切模 $I_T/I_T^2$ 的 Fitting ideal 由一个正则序列控制，
所以
\[
  \operatorname{length}_{\Z_\ell}(\Z_\ell/\Phi_T)
"]=
  \operatorname{length}_{\Z_\ell}(I_T/I_T^2)
\]
恰好达到由生成元数与关系数给出的最小值。
上面的同构已经说明 $I_R/I_R^2\cong I_T/I_T^2$，
故 $R$ 有与 $T$ 相同的变量数与关系数。
de Smit--Lenstra 与 Wiles 的交换代数判据于是推出 $R$ 也必须是 complete intersection。
在 complete intersection 范畴里，congruence ideal 与 cotangent 长度完全相配，
所以不等式自动变成等式，最终得到
\[
  R\cong T.
\]
\end{proof}

Taylor--Wiles 的决定性发明是引入一族辅助素数 $Q_n$，
构造带有额外级别结构的环 $R_{Q_n}$ 与 $T_{Q_n}$，
再通过 inverse limit 与 patching 把问题抬升到一个更自由的极限情形，
在极限上先证明
\[
  R_\infty\cong T_\infty,
\]
然后再下降回原始的最小情形。

\subsection*{补充：Selmer 群如何进入 $R=T$}
CSS 对 Wiles 方法的改编特别强调：
所谓“比较两个局部环大小”，
真正承担数值信息的对象并不是环本身的显式方程，
而是由 Galois 上同调定义出的 Selmer 群。

\begin{definition}[上同调定义的 Selmer 群]
设 $M$ 是一个带 $G_\Q$ 作用的离散 $\ell$-初群，
并在每个素数 $v$ 固定一个局部条件
\[
  L_v\subset H^1(\Q_v,M).
\]
则对应的 Selmer 群定义为
\[
  \Sel_{\mathcal L}(\Q,M)
  =\ker\!\left(H^1(\Q,M)\to \prod_v H^1(\Q_v,M)/L_v\right).
\]
对椭圆曲线 $E/\Q$ 与 $M=E[\ell]$，
通常取 $L_v$ 为 Kummer 映射像，得到经典的 $\ell$-Selmer 群。
\end{definition}

\begin{proposition}[Selmer 群与 Tate--Shafarevich 群]
\label{prop:ch10-selmer-sha}
对椭圆曲线 $E/\Q$ 与素数 $\ell$，存在正合列
\[
  0\to E(\Q)\otimes \Z/\ell\Z
  \to \Sel(\Q,E[\ell])
  \to \text{Tate--Shafarevich 群的 }\ell\text{-挠部分}
  \to 0.
\]
因此 Selmer 群同时测量了有理点秩与 Tate--Shafarevich 群的失败局部--整体原理。
\end{proposition}

\begin{proof}[说明]
Kummer 正合列
\[
  0\to E[\ell]\to E(\bar\Q)\xrightarrow{\ell} E(\bar\Q)\to 0
\]
在全局与局部 Galois 上同调之间给出连接映射。
把所有局部可解条件同时施加到 $H^1(\Q,E[\ell])$ 上，
便得到 Selmer 群；
其商正好记录那些在每个 $\Q_v$ 上都可解、但不来自全局点的扭类，
也就是 Tate--Shafarevich 群的相应部分。
\end{proof}

\begin{theorem}[Wiles 数值判据中的对偶 Selmer 公式]
\label{thm:ch10-wiles-selmer-duality}
设 $f$ 是与极大理想 $\mathfrak m$ 对应的模形式，$T_f$ 是其 $\ell$-进 Galois 晶格，$V_f=T_f\otimes \Q_\ell$。
则在适当的最小局部条件下，Poitou--Tate 对偶给出
\[
  \bigl|\Sel(\Q,V_f/T_f)\bigr|
  =
  \bigl|\Sel(\Q,T_f^*(1)/V_f^*(1))\bigr|,
\]
这里右边是对偶表示的 Selmer 群。
更常见地，人们把它写成两边长度或余维相等。
\end{theorem}

\begin{proof}[说明]
局部 Tate 对偶把每个 $H^1(\Q_v,-)$ 上的允许局部条件与其正交补配对起来；
Poitou--Tate 全局对偶则把这些局部配对拼成一个全局 Euler 特征公式。
在 Wiles 所需的最小情形里，
局部条件被精心选择到恰好自对偶，
于是原始 Selmer 群与对偶 Selmer 群拥有相同的有限大小。
这就是切空间长度可以被精确配平的根本原因。
\end{proof}

\textbf{注。}
形变环 $R$ 的 cotangent space 大小由原始 Selmer 群控制，
而 Hecke 环 $T$ 的 congruence ideal 则由模形式间的同余控制。
Wiles 的数值判据所做的，正是借助上面的对偶性把两边的长度比较变成一个闭合的算术恒等式；
这也解释了为什么 Selmer 群是 $R=T$ 定理里的真正“计量单位”。

\begin{theorem}[Ribet 水平下降]
\label{thm:ribet}
设 $p\nmid N$，
$f\in \Sk_2(\GammaO(Np))$ 是归一化 $p$-new 特征新形式，
$\ell\neq p$，且 residual 表示
\[
  \rhobar_{f,\ell}\colon G_\Q\to \GL_2(\F_\ell)
\]
绝对不可约。
若 $\rhobar_{f,\ell}$ 在 $p$ 处无分歧，
则存在归一化新形式
\[
  g\in \Sk_2(\GammaO(N))
\]
使得
\[
  \rhobar_{g,\ell}\cong \rhobar_{f,\ell}.
\]
\end{theorem}

\begin{proof}[证明思路]
这是第\ref{ch:galois}章的核心定理之一。
其本质是：
级别 $Np$ 与级别 $N$ 的模曲线之间存在两条退化映射，
它们生成 $p$-old 部分；
若 residual 表示在 $p$ 处已经无分歧，
则它在 $J_0(Np)[\mathfrak m]$ 中留下的痕迹无法纯粹待在真正的 $p$-new 部分里，
最终会被 Ihara 引理逼回 old 部分。
而 old 部分本来就来自级别 $N$，
故同一份 residual 表示也来自级别 $N$ 的新形式。
\end{proof}

\subsection*{第五步：$3$--$5$ 技巧处理可约例外}
Wiles 证明里最著名的巧思之一，
就是当 $\rhobar_{E,3}$ 不可用时，
并不是随便换一个素数，
而是精确地换到 $5$。

\begin{proposition}[$3$--$5$ 技巧]
\label{prop:ch10-35-trick}
设 $E/\Q$ 是半稳定椭圆曲线。
若 $\rhobar_{E,3}$ 可约，
则 $\rhobar_{E,5}$ 绝对不可约，且 $\rhobar_{E,5}$ 是模的。
更精确地，存在另一条半稳定椭圆曲线 $A/\Q$，满足
\[
  A[5]\cong E[5],
  \qquad
  \rhobar_{A,3}\ \text{绝对不可约}.
\]
于是 $A$ 模，进而 $\rhobar_{E,5}$ 模；再由提升定理，$E$ 亦模。
\end{proposition}

\begin{proof}
这条技巧的核心不是直接处理 $E[5]$，
而是先找一条与 $E$ 共享 $5$-挠 Galois 模、但在 $3$ 上更“好用”的新曲线 $A$。
整个论证分三步。

\textbf{第一步：$\rhobar_{E,3}$ 与 $\rhobar_{E,5}$ 不能同时可约。}
若 $\rhobar_{E,3}$ 可约，
则 $E[3]$ 中存在一个 $G_\Q$-稳定的一维子空间，
也就是一个有理循环子群 $C_3\subset E[3]$，故 $E$ 有一个定义在 $\Q$ 上的 $3$-isogeny。
同理，若 $\rhobar_{E,5}$ 也可约，
则存在有理循环子群 $C_5\subset E[5]$，故 $E$ 有一个 $5$-isogeny。
因为 $3$ 与 $5$ 互素，
群 $C_3\oplus C_5\subset E[15]$ 是 $G_\Q$-稳定的，且同构于循环群 $C_{15}$；
因此 $E$ 给出模曲线 $X_0(15)$ 上的一个有理点，也就是说 $E$ 有一个有理 $15$-isogeny。

但 $X_0(15)(\Q)$ 的非尖点早已被完全列出；
对应的椭圆曲线都在 $2$ 或 $5$ 处具有加性约化，因而都不是半稳定的。
这与 $E$ 半稳定矛盾。
故在半稳定情形里，$\rhobar_{E,3}$ 可约必然推出 $\rhobar_{E,5}$ 不可约；
而二维 mod $5$ 表示若可约于某个扩域上，
则其中一条 Jordan--H\"older 线在该扩域上稳定，
再由 Galois 共轭可知原表示已有一条稳定直线，
所以这里实际上得到的是绝对不可约。

\textbf{第二步：构造一条与 $E[5]$ 相同、但 $A[3]$ 不可约的曲线。}
考虑模曲线 $X_E(5)$，它参数化的是带有 symplectic 同构
\[
  \alpha\colon B[5]\xrightarrow{\sim} E[5]
\]
的椭圆曲线 $B/\Q$。
因为 $E$ 自己给出一个有理点，$X_E(5)$ 是有理曲线；
而 $X(5)$ 的亏格为 $0$，故 $X_E(5)\cong \mathbb P^1_\Q$，因此有无穷多个有理点。

现在考察其中使 $B[3]$ 可约的点。
“$B[3]$ 可约”意味着 $B$ 在模曲线 $X_0(3)$ 上有有理像；
 于是这些点落在从 $X_E(5)$ 到 $X(15)$ 的一个薄子集里。
由于 $X_E(5)\cong \mathbb P^1$，Hilbert 不可约定理告诉我们：
在无穷多个有理点中，可以选到一点对应的椭圆曲线 $A$，使得
\[
  A[5]\cong E[5],
  \qquad
  A[3]\ \text{绝对不可约}.
\]
进一步通过在有限多个局部点上施加开条件，还可把 $A$ 选成半稳定的。
这就是 $3$--$5$ 技巧最精细的几何输入。

\textbf{第三步：把 $A$ 的模性传回 $E$。}
由于 $A[3]$ 绝对不可约，且其像仍落在 $\GL_2(\F_3)$ 中，
\cref{lem:ch10-langlands-tunnell} 适用，故 $\rhobar_{A,3}$ 是模的。
再对 $A$ 应用 Wiles 的提升定理，得到 $A$ 本身是模的。
于是 $A$ 的 $5$ 进表示模，故其残余表示 $A[5]$ 也模。
由构造 $A[5]\cong E[5]$，可知
\[
  \rhobar_{E,5}\cong \rhobar_{A,5}
\]
也是模的。
最后对 $E$ 的 $5$ 进表示应用同一个模性提升定理，
便得到 $E$ 模。

因此，当 $\rhobar_{E,3}$ 可约时，
我们先用另一条曲线 $A$ 把模性搬运到 $\rhobar_{E,5}$，
再由 $5$ 进提升回到 $E$；
这就完成了 $3$--$5$ 技巧的论证。
\end{proof}

\textbf{注。}
从今天的视角看，
“先证残余模性，再证所有允许的提升都是模的”
已经是一个标准模板。
但在 Wiles 之前，
真正新颖的正是把这个模板组织成可计算的局部环比较问题。
下一节我们就专门解释这个比较问题里的技术词：形变函子、形变环与 $R=T$。
